
\subsection{Qualitative Case Study: Resolution of Modality Conflict}
\label{subsec:case_study}
To demonstrate the framework's internal decision-making process under discordance, we trace the inference trajectory of a query exhibiting a deliberate semantic conflict.

\textbf{1. Input Processing:}
The system receives a composite query with conflicting clinical signals:
\begin{itemize}
    \item \textbf{Text ($Q_{txt}$):} ``Hi, I hit the side of my forehead above my eyebrow very hard 2 days ago. It was only a lump at first then the lump went down and the side of my forehead swelled up. Been having head pain and a lot of tingles in my head almost like it itches and now my eyelid is black and swollen...'' (Trauma/Hematoma narrative).
    \item \textbf{Image ($Q_{img}$):} A dermoscopic image of a \textit{Skin Growth} (seborrheic keratosis).
\end{itemize}

\textbf{2. Constraint Verification:}
Before generation, the Constraint Validator evaluates the retrieved set:
\begin{itemize}
    \item \textit{Modality Consistency Check:} The cosine similarity between $Q_{txt}$ and $Q_{img}$ is computed as $\mathbf{0.2685}$ (Low).
    \item \textit{Action:} This violation explicitly flags the query as \textit{"Ambiguous/Conflicting"}, triggering the counterfactual sweep.
\end{itemize}

\textbf{3. Fusion and Initial Retrieval:}
The Adaptive Fusion module computes a joint embedding. Due to the high visual feature clarity, the fusion weights heavily incentivize the visual signal, but the strong textual narrative pulls the distribution.
\begin{itemize}
    \item \textit{Retrieved Hypotheses:} The system retrieves two competing concept clusters: \textit{Skin Growth} ($Score=0.51$) and \textit{Swollen Eye} ($Score=0.21$).
    \item \textit{Initial State:} A conventional RAG system would naively output \textit{Skin Growth} as the top-1 prediction.
\end{itemize}

\textbf{4. Counterfactual Disentanglement:}
The system isolates the signals to test hypothesis stability:
\begin{enumerate}
    \item \textit{Text-Neutralized (Visual-only):} The distribution collapses to \textit{Skin Growth} ($p=1.0$), confirming the image is valid but unrelated to the text.
    \item \textit{Image-Neutralized (Text-only):} The distribution shifts entirely to \textit{Swollen Eye} ($p=0.85$), confirming the text describes a different pathology.
\end{enumerate}

\textbf{5. Final Validated Outcome:}
Instead of hallucinating a forced reconciliation, CountRAG-Clinic outputs a \textbf{Low Robustness} warning (Stability Divergence $= 0.39$) and presents both diagnoses with their respective un-mixed evidence sources. This "Safe Failure" prevents the model from ignoring the patient's description of trauma in favor of the visual finding, prompting clinician review.

\textbf{6. System Output (API Response):}
\begin{verbatim}
{
  "stability": {
    "robustness": "low", 
    "divergence": 0.39 
  },
  "constraints": { 
    "modality_consistency": 0.26, 
    "status": "VIOLATION" 
  },
  "outcome": "AMBIGUOUS_CONFLICT"
}
\end{verbatim}

\textit{``The visual evidence suggests a skin growth (e.g., keratosis), but the patient's description of recent head trauma and swelling is consistent with a hematoma ('Swollen Eye'). Due to this conflict (Modality Consistency: 0.26), both possibilities are presented. Clinical verification is recommended to rule out trauma complications."}
